\section{A third-indicator triangle obstruction at conductor forty}
\label{sec:q630-weighted-indicator}

\begin{theorem}\label{thm:q630-exclusion}
The fusion ring $R8\_2\_630$ has no complex categorification.
\end{theorem}
\begin{proof}
Proposition~\ref{prop:hi6-pair-galois-pu} reduces a hypothetical complex
categorification to a pseudo-unitary Galois conjugate. Work with its
canonical positive spherical structure. Put $s=\sqrt{10}$ and
$K=\mathbb Q(s)$. In census order, the FP dimensions and global
dimension are
\[
(d_0,\ldots,d_7)=(1,1,2+s,3+s,3+s,2+s,4+s,4+s),
\qquad D=120+36s.
\]
The complete character table gives formal codegrees
\[
D,\quad D'=120-36s,\quad 4,6,6,12,12,12.
\]
The source ring is commutative, so each constituent $A_E$ of $I(1)$
has multiplicity one and dimension $D/f_E$
\cite[Theorem~2.13]{OstrikPiv3}. Let $J$ denote the constituent for
$D'$; then $\dim J=D/D'=19+6s$.

The multiplicity vector of $I(1)$ satisfies $Sm=m$, so its support
is Galois-stable with positive signs. The central unit has Galois
orbit exactly $\{1,J\}$: these are the unique constituents of
$I(1)$ of dimensions $1$ and $D/D'$. Thus $w=e_1-e_J$ has the
primitive quadratic torus character $\chi_{40}$. The known twist-one
block satisfies
\[
\operatorname{proj}_{I(1)} Sw
 =\left(\frac1D-\frac1{D'}\right)w
 =-\frac{s}{20}w=-\frac{w}{\sqrt{40}}.
\]
Indeed the $1,J$ submatrix has diagonal $1/D$ and off-diagonal
$1/D'$, while all other entries in the two columns agree since their
codegrees are rational.

The local compression bound \cite[Lemma~4.1]{SchDimensionBlock} is
therefore attained. Its equality refinement, proved in the argument
of Lemma~\ref{lem:q17-conductor-equality}, works at each prime power:
for conductor $p^c$, the equality vectors are pullbacks from level
$p^c$; trivial local characters force the trivial representation.
Applying this to the nontrivial local conductors $8$ and $5$ gives
level $40$ for the cyclic representation generated by $w$.
The general equality proof is recorded in
\path{compression/CONDUCTOR_EQUALITY_LEMMA.md}.
Since the coordinate of $Sw$ at a center simple of dimension $z$ is
\[
\frac zD-\sigma\!\left(\frac zD\right),
\]
where $\sigma$ flips $K$, every such object with $z/D$ nonrational
has $\theta^{40}=1$.

Apply the same saturated Galois indicator lemma used for the five
conductor-seventeen cases, Lemma~\ref{lem:saturated-galois-indicator},
now with
\[
\ell(a+bs)=3a-10b,\qquad f=40,\qquad k=3.
\]
We have $\ell(D)=0$, and $\chi_{40}(3)=1$ means that
$\zeta_{40}\mapsto\zeta_{40}^3$ fixes $K$. The source weights are
\[
(\ell(d_i))_{i=0}^7=(3,3,-4,-1,-1,-4,2,2),
\]
so the lemma gives
\begin{equation}\label{eq:q630-indicator-balance}
\sum_{i\ne0}\ell(d_i)\nu_3(X_i)=0.
\end{equation}
This is an exact complex identity. The first trace is transformed
into the third by a field automorphism fixing every coefficient;
no reality of the third indicators is assumed.

Finally the source tensor gives
\[
\bigl(\dim\operatorname{Hom}(1,X_i^3)\bigr)_{i=0}^7
 =(1,0,0,1,1,1,0,0).
\]
The only nonzero nonunit third indicators are consequently those
at labels $3,4,5$. Each is the trace of an order-three rotation on
a one-dimensional space, hence a cube root of unity of modulus one.
Equation~\eqref{eq:q630-indicator-balance} becomes
\[
\nu_3(X_3)+\nu_3(X_4)+4\nu_3(X_5)=0.
\]
The first two terms have sum of absolute value at most $2$, whereas
the last has absolute value $4$. This strict triangle contradiction
excludes the hypothetical category and proves the theorem.
\end{proof}

The source-bound standard-library verifier
\path{verify_r630_weighted_indicator.py} replays the structural
sphericality gate, complete character table, dimensions, codegrees,
conductor arithmetic, weights, and cubic invariant dimensions.
The theoretical contradiction uses no enumeration of center rows.
Its normal and optimized receipts and independent audit accompany
\path{R630_WEIGHTED_GALOIS_INDICATOR_PROOF.md} and the common lemma
\path{SATURATED_GALOIS_INDICATOR_LEMMA.md} in
\path{evidence/unified_20261016/rank8/}.
